\section{Introduction}
\label{sec:introduction}

Elliptic-curve cryptography routinely transmits points on a fixed
public curve, and such points are easily distinguished from uniform
random bit strings: an observer who can tell elliptic-curve traffic
apart from random traffic can selectively block it. Bernstein,
Hamburg, Krasnova and Lange's Elligator~\cite{bernstein2013a}
addresses this by giving efficient, invertible encodings between
(most of) the points of a suitable elliptic curve and uniform bit
strings, so that a random curve point can be transmitted as what
looks like a random string. The paper introduces two such encodings,
Elligator~1 and Elligator~2, together with Curve1174, a concrete
Edwards curve suited to Elligator~1. Later work has extended the
approach to further curve shapes, including a Legendre-curve variant
due to Saha and Karati~\cite{gourab2026a}.

Despite substantial practical uptake --- Elligator has been
numerically tested~\cite{elligatorcryptotests} and implemented
across a range of libraries and protocols~\cite{elligatororg} --- the
underlying theory has, to our knowledge, never been formalized in a
proof assistant. This paper reports on such a formalization, carried
out in CSLib, a Lean~4 library that currently has essentially no
cryptographic primitives, not even basic elliptic-curve definitions.
Elligator is accordingly the first substantial elliptic-curve
encoding family to enter CSLib.

\paragraph{Contributions.}
\begin{enumerate}
  \item A formalization of elliptic curves and the curve primitives
    needed to state Elligator, none of which previously existed in
    CSLib.
  \item A formalization of Elligator~1 and Curve1174, and a
    generalization of the viability lemmas underlying Elligator~1
    from prime fields to prime power fields (Section~\ref{sec:elligator1}).
  \item A formalization of Elligator~2 and of the Legendre-curve
    variant of Saha and Karati as instances of a common encoding
    family alongside Elligator~1 (Section~\ref{sec:elligator2-legendre}).
  \item Computable, extractable definitions throughout, so that the
    formalized encodings are not merely existence statements but can
    be run and checked against published test
    vectors (Section~\ref{sec:computability}).
\end{enumerate}

\paragraph{Outline.}
Section~\ref{sec:background} recalls the curve and field background.
Section~\ref{sec:elligator1} covers Elligator~1 and its
generalization. Section~\ref{sec:elligator2-legendre} covers
Elligator~2 and the Legendre-curve variant.
Section~\ref{sec:curve1174} covers Curve1174.
Section~\ref{sec:computability} discusses computability and
validation against test vectors. Section~\ref{sec:related-work}
discusses related work and Section~\ref{sec:conclusion} concludes.

